% Capítulo — TCC 1
\chapter{Fundamentação teórica}
\label{chap:fundamentacao}

% Abertura do capítulo: um parágrafo dizendo o que ele apresenta.

\section{Integração de dados e schema matching}
\label{sec:fundamentacao-integracao-de-dados-e-schema-matching}
% Lenzerini (2002); Doan, Halevy e Ives (2012); Rahm e Bernstein (2001).

\section{Perfilamento de dados e evidências de compatibilidade}
\label{sec:fundamentacao-perfilamento-de-dados-e-evidencias-de-compatibilidade}
% Abedjan, Golab e Naumann (2015).

\section{Modelos de linguagem}
\label{sec:fundamentacao-modelos-de-linguagem}
% Saída estruturada, DSPy, cadeia de raciocínio (Khattab et al., 2024; Wei et al., 2022).

\section{Arquitetura de dados}
\label{sec:fundamentacao-arquitetura-de-dados}
% Camadas bronze, silver e gold; orquestração orientada a dados (Armbrust et al., 2021; Reis e Housley, 2022).

\section{Spec-Driven Development e agentes de IA de código}
\label{sec:fundamentacao-spec-driven-development-e-agentes-de-ia-de-codigo}
% Böckeler (2025); GitHub Spec Kit; Meyer (1992); Adzic (2011). Literatura acadêmica ainda escassa: buscar ICSE, FSE, arXiv.

\section{Dados abertos governamentais no Brasil}
\label{sec:fundamentacao-dados-abertos-governamentais-no-brasil}
% LAI, Decreto 8.777/2016, Lei 14.129/2021, LGPD.

\section{Artigo-base: SPAPI-Tester}
\label{sec:fundamentacao-artigo-base-spapi-tester}
% Problema, arquitetura, categorias de atrito e resultados; o que o TCC adota.
